% !TeX root = ../main_skripsi.tex
\chapter{KAJIAN PUSTAKA}

\section{Kajian Teori}

[Sintesis definisi, konsep, asumsi, dan indikator tiap variabel dari sumber
primer atau sumber ilmiah kredibel. Hindari rangkaian kutipan tanpa analisis dan
jangan memakai materi pembelajaran yang belum melalui uji publik sebagai dasar
utama.]

\section{Hasil Penelitian yang Relevan}

[Bandingkan penelitian terdahulu secara naratif untuk menunjukkan persamaan,
perbedaan, keterbatasan, celah pengetahuan, dan posisi penelitian ini.]

% Contoh struktur tabel. Hapus atau isi hanya memakai paper terverifikasi.
\begin{table}[H]
  \centering
  \caption{Matriks hasil penelitian yang relevan}
  \label{tab:penelitian-relevan}
  \begin{tblr}{width=\textwidth,colspec={Q[c,0.07\textwidth] X X X}}
    No. & Penulis (Tahun) & Metode dan Temuan & Perbedaan dengan Penelitian \\
    1 & [Belum diisi] & [Belum diisi] & [Belum diisi] \\
  \end{tblr}
\end{table}

\section{Kerangka Berpikir}

[Susun hubungan logis antarvariabel berdasarkan sintesis teori dan bukti
penelitian terdahulu. Tambahkan diagram jika alur konseptual sudah tervalidasi.]

\section{Pertanyaan Penelitian dan/atau Hipotesis}

[Tulis pertanyaan penelitian atau hipotesis yang merupakan jabaran langsung
dari rumusan masalah. Gunakan hipotesis hanya bila desain memang mengujinya.]
